\section{Ablation}
\label{sec:ablation}

Each component of the proposed monitors is removed or replaced in turn on the same generators, with one change at a time and the same streams for every variant of a scenario, so that the differences within a table are paired and smaller than the differences between runs. The three tables below are the ablation: \cref{tab:ablation-grid} for the mixture of the rate monitor, \cref{tab:ablation-battery} for the combination of the two VaR monitors,
and \cref{tab:ablation-pd} for the PD monitor.

\subsection{The mixture over alternative hit rates}
\label{sec:ablation-mixture}

\Cref{tab:ablation-grid} replaces the full grid by single alternatives and by parts of the grid. A single alternative is the likelihood-ratio e-process of one hit rate and is the monitor that the theory says is fastest when its alternative is the truth. At a hit rate three times the nominal, the single alternative at $3\alpha$ has a median of \nAbMedThreeSingleThree\ days and \nAbByThreeSingleThree\,\% power by day 250,
and at four times the nominal the one at $4\alpha$ has \nAbMedFourSingleFour\ days and \nAbByFourSingleFour\,\%; against them the full grid has \nAbMedThreeFull\ days and \nAbByThreeFull\,\% at three times and \nAbMedFourFull\ days and \nAbByFourFull\,\% at four times. The difference is the price of not knowing the severity. A single alternative that is far from the truth is worse than the mixture: the alternative at $10\alpha$ has a median beyond the horizon at three times
the nominal (\nAbMedThreeSingleTen) and only \nAbByThreeSingleTen\,\% power by day 250, and the alternative at $1.5\alpha$ has \nAbByThreeSingleOneFive\,\% there.

Cutting the grid shows which alternatives carry the power. Dropping the two alternatives below the nominal level (the upper half) raises the power by day 250 at three times the nominal from \nAbByThreeFull\ to \nAbByThreeUpper\,\% and lowers the median from \nAbMedThreeFull\ to \nAbMedThreeUpper\ days at no visible cost in size (\nAbNullFull\ and \nAbNullUpper\,\%), because the weight on the alternatives that point the other way is wasted when the departure is upward.
Keeping only the alternatives $1.5$ to $4$ gives \nAbByThreeUpperLow\,\% at three times the nominal and, among the cuts of the grid, the best power against twice the nominal (\nAbEverTwoUpperLow\,\% rejected within 1,000 days), while keeping only $3$ to $10$ gives the fastest delays at four and six times and loses power at twice the nominal (\nAbEverTwoUpperHigh\,\%). The grid trades breadth of coverage against speed at a given severity, and no cut dominates; the full grid is the choice that does not prejudge the severity or the direction.

\begin{table*}[htbp]
\centering
\small
\caption{Ablation of the mixture over alternative hit rates of the rate monitor on 2,000 streams of 1,000 days per scenario: the full grid, its upper half, the upper half cut at its low and high ends, and each single alternative, which is the likelihood-ratio e-process of one hit rate. Size and power (x2) are shares ever rejected; power by day 250 and median days are for hit rates three (x3), four (x4) and six (x6) times too high; brackets are 95\,\% intervals.}
\label{tab:ablation-grid}
\resizebox{\linewidth}{!}{%
\begin{tabular}{p{4.4cm}rrrrrrr}
\toprule
Mixture grid & Size, \% & Power x2, \% & Power by 250 x3, \% & Median day x3 & Power by 250 x4, \% & Median day x4 & Median day x6 \\
\midrule
\relax Full two-sided grid, $0.25$ to $10$ & \relax 2.2 [1.7, 3.0] & \relax 55.8 [53.6, 58.0] & \relax 52.3 [50.1, 54.5] & \relax 232 [220, 249] & \relax 77.8 [75.9, 79.6] & \relax 128 [121, 133] & \relax 58 [56, 60] \\
\relax Upper half, $1.5$ to $10$ & \relax 2.4 [1.8, 3.2] & \relax 59.2 [57.1, 61.4] & \relax 57.1 [54.9, 59.3] & \relax 215 [194, 226] & \relax 82.0 [80.2, 83.6] & \relax 108 [105, 120] & \relax 56 [54, 58] \\
\relax Upper alternatives $1.5$ to $4$ & \relax 2.6 [2.0, 3.4] & \relax 64.4 [62.3, 66.5] & \relax 57.4 [55.2, 59.6] & \relax 207 [199, 214] & \relax 82.5 [80.8, 84.2] & \relax 120 [110, 130] & \relax 60 [59, 62] \\
\relax Upper alternatives $3$ to $10$ & \relax 3.0 [2.4, 3.9] & \relax 53.4 [51.3, 55.6] & \relax 59.3 [57.1, 61.4] & \relax 191 [180, 203] & \relax 83.1 [81.4, 84.7] & \relax 104 [99, 107] & \relax 50 [42, 53] \\
\relax Single alternative $1.5\alpha$ & \relax 0.9 [0.5, 1.4] & \relax 59.5 [57.3, 61.6] & \relax 18.2 [16.6, 20.0] & \relax 381 [374, 406] & \relax 46.6 [44.4, 48.7] & \relax 261 [255, 267] & \relax 154 [150, 158] \\
\relax Single alternative $2\alpha$ & \relax 2.9 [2.3, 3.7] & \relax 72.5 [70.6, 74.5] & \relax 51.9 [49.7, 54.1] & \relax 245 [239, 251] & \relax 78.7 [76.9, 80.4] & \relax 162 [158, 167] & \relax 95 [93, 96] \\
\relax Single alternative $3\alpha$ & \relax 3.3 [2.6, 4.2] & \relax 63.3 [61.2, 65.4] & \relax 64.0 [61.9, 66.1] & \relax 175 [170, 181] & \relax 86.4 [84.8, 87.8] & \relax 110 [106, 114] & \relax 60 [59, 62] \\
\relax Single alternative $4\alpha$ & \relax 3.0 [2.4, 3.9] & \relax 48.1 [46.0, 50.3] & \relax 63.0 [60.9, 65.1] & \relax 167 [159, 173] & \relax 85.2 [83.6, 86.7] & \relax 99 [91, 105] & \relax 55 [53, 56] \\
\relax Single alternative $10\alpha$ & \relax 2.0 [1.5, 2.7] & \relax 14.4 [12.9, 16.0] & \relax 38.5 [36.4, 40.7] & \relax $>$1,000 & \relax 63.9 [61.8, 66.0] & \relax 117 [108, 135] & \relax 42 [40, 43] \\
\bottomrule
\end{tabular}}
\end{table*}


\subsection{The direction of the alternatives}
\label{sec:ablation-direction}

The upper half of the grid is the one-sided monitor of a supervisor, who penalises a VaR that understates the risk and does not penalise one that is too conservative. It is the cheapest change in \cref{tab:ablation-grid}, and it is the one that moves the hypotheses. On the pooled \nSeedsPooledN\ streams per factor of \cref{tab:x3-seeds} the upper-side monitor has a median of \nSeedsUpThreeMed\ days at three times the nominal level against \nSeedsThreeMed\ for the two-sided one, and a power cost of \nSeedsUpThreeCost\ points against \nSeedsThreeCost;
at four times the nominal, \nSeedsUpFourMed\ days against \nSeedsFourMed\ and \nSeedsUpFourCost\ points against \nSeedsFourCost. The upper-side monitor rejects a correct model in \nNullSizeEUp\,\% of the main-run series, against \nNullSizeE\,\% for the two-sided one: its size is not larger in any way that matters. The price is that it has no power against a model that is too conservative, which the supervisory test does not ask for and a validator of the efficiency of a model may; for that validator the two-sided monitor is the choice, and its delay at a hit rate below the nominal is not measured here.

\subsection{The combination of the two VaR monitors}
\label{sec:ablation-battery}

\Cref{tab:ablation-battery} runs the rate monitor, the clustering monitor, their average, the union of the two at a threshold of 40 each (the Bonferroni split of the size, valid by the union bound) and their product on the same streams. The average and the union are valid and almost identical: they reject \nAbBatThreeAverage\ and \nAbBatThreeUnion\,\% of the streams with a hit rate three times too high, against \nAbBatThreeRate\,\% for the rate monitor alone, and \nAbBatClFourAverage\ and \nAbBatClFourUnion\,\% of the streams with clusters at $q=0.4$, against
\nAbBatClFourCluster\,\% for the clustering monitor alone. Averaging two monitors costs at most the factor two on each, and it buys the detection of either departure, which neither monitor has alone. The product is not an e-process for two monitors of one stream (\cref{prop:product}), but its measured size at a level of 1\,\% is \nAbBatNullProduct\,\%: the rate monitor and the clustering monitor rarely move together on a correct model at that level, so the invalidity does not show at 1\,\% and does show at the level of
$0.2$ of the enumeration (\cref{tab:null-expectation}). The product detects \nAbBatClFourProduct\,\% at $q=0.4$, about the same as the average (\nAbBatClFourAverage\,\%), so nothing is gained by taking the risk.

\begin{table*}[htbp]
\centering
\small
\caption{Ablation of the combination of the two VaR monitors that read the same stream, on 2,000 streams of 1,000 days per scenario: share of streams ever rejected, in percent, when the model is right (size), when the hit rate is three times too high, and when hits cluster at the right rate. The product is shown to measure its size, which the theory does not bound.}
\label{tab:ablation-battery}
\resizebox{\linewidth}{!}{%
\begin{tabular}{p{5.4cm}rrrrr}
\toprule
Combination & Size, \% & Rate x3, \% & Clustering $q=0.2$, \% & Clustering $q=0.4$, \% & Median day, $q=0.4$ \\
\midrule
\relax Rate e-process & \relax 1.8 [1.3, 2.5] & \relax 97.5 [96.8, 98.1] & \relax 7.5 [6.4, 8.7] & \relax 18.1 [16.5, 19.9] & \relax $>$1,000 \\
\relax Clustering e-process & \relax 0.4 [0.2, 0.9] & \relax 8.6 [7.4, 9.9] & \relax 53.3 [51.2, 55.5] & \relax 77.2 [75.4, 79.0] & \relax 525 [491, 556] \\
\relax Average of the two (battery, valid) & \relax 1.2 [0.8, 1.8] & \relax 96.4 [95.5, 97.1] & \relax 51.6 [49.5, 53.8] & \relax 78.8 [77.0, 80.6] & \relax 523 [487, 554] \\
\relax Either at threshold 40 (union, valid) & \relax 1.2 [0.8, 1.8] & \relax 96.4 [95.4, 97.1] & \relax 51.6 [49.5, 53.8] & \relax 78.8 [77.0, 80.5] & \relax 525 [491, 557] \\
\relax Product of the two (not an e-process on one stream) & \relax 2.5 [1.9, 3.2] & \relax 95.5 [94.4, 96.3] & \relax 48.8 [46.6, 50.9] & \relax 79.0 [77.1, 80.7] & \relax 474 [445, 506] \\
\bottomrule
\end{tabular}}
\end{table*}


\subsection{The PD monitor: combination, direction and grid}
\label{sec:ablation-pd}

\Cref{tab:ablation-pd} varies the combination of the grades, the direction and the grid of the PD monitor. With the PD right, none of the monitors rejects more than \nPdOneSingle\,\% of the histories; the chi-square test read once at year 10 is at \nPdOneChiOne\,\% and read every year at \nPdOneChiEvery\,\%. With every PD multiplied by $1.1$, the combination by product rejects \nPdOneOneProd\,\%, by average \nPdOneOneAvg\,\%, and the upper-side average \nPdOneOneAvgUp\,\% of the histories, so the direction and the combination matter little against the grid. The grid matters a great deal: with four alternatives between $1.1$ and $2$ instead of the
sixteen of the default, the upper-side average rejects \nPdOneOneCoarse\,\% of the histories at $1.1$ and \nPdOneZeroFiveCoarse\,\% at $1.05$, against \nPdOneOneAvgUp\ and \nPdOneZeroFiveAvgUp\,\%, with a size of \nPdOneCoarse\,\%. A single alternative at $1.25$ gives \nPdOneOneSingle\,\% at $1.1$. The default grid spreads the weight over alternatives far from the departures that a small relative error produces: half of its upper alternatives are above $2$, and
\cref{prop:growth} says that a mixture pays the logarithm of the number of alternatives it does not need. The coarse grid still detects less than the chi-square rule at a nominal size of 5\,\% (\nPdcOneOneChiSim\,\% at $1.1$ and \nPdcOneZeroFiveChiSim\,\% at $1.05$; \cref{tab:pd-equal-size}). A validator who knows the range of miscalibrations of interest can set a coarse grid before the data and gain power at no cost in size; the default of the engine was fixed before these experiments and is not changed here.

\begin{table*}[htbp]
\centering
\small
\caption{Ablation of the PD monitor of five grades over ten years on 1,000 histories per column from the generator of default counts with no asset correlation: share of histories ever rejected, in percent with 95\,\% Wilson intervals, for the grades combined by product or by average, two-sided or upper, with coarser grids, and for the chi-square test across grades read once or every year.}
\label{tab:ablation-pd}
\resizebox{\linewidth}{!}{%
\begin{tabular}{p{5.8cm}rrrrr}
\toprule
Rule & Every PD right (size) & Every PD $\times 1.05$ & Every PD $\times 1.1$ & Every PD $\times 1.25$ & 10\,\% grade $\times 1.3$ \\
\midrule
\relax Product, two-sided & \relax 0.0 [0.0, 0.4] & \relax 0.0 [0.0, 0.4] & \relax 26.7 [24.1, 29.5] & \relax 100.0 [99.6, 100.0] & \relax 100.0 [99.6, 100.0] \\
\relax Product, upper & \relax 0.0 [0.0, 0.4] & \relax 0.0 [0.0, 0.4] & \relax 28.4 [25.7, 31.3] & \relax 100.0 [99.6, 100.0] & \relax 99.9 [99.4, 100.0] \\
\relax Average, two-sided & \relax 0.4 [0.2, 1.0] & \relax 2.9 [2.0, 4.1] & \relax 55.9 [52.8, 58.9] & \relax 100.0 [99.6, 100.0] & \relax 100.0 [99.6, 100.0] \\
\relax Average, upper & \relax 0.2 [0.1, 0.7] & \relax 4.7 [3.6, 6.2] & \relax 60.0 [56.9, 63.0] & \relax 100.0 [99.6, 100.0] & \relax 100.0 [99.6, 100.0] \\
\relax Average, upper, four alternatives $1.1$ to $2$ & \relax 0.7 [0.3, 1.4] & \relax 19.9 [17.5, 22.5] & \relax 94.5 [92.9, 95.8] & \relax 100.0 [99.6, 100.0] & \relax 100.0 [99.6, 100.0] \\
\relax Average, upper, one alternative $1.25$ & \relax 0.9 [0.5, 1.7] & \relax 9.7 [8.0, 11.7] & \relax 60.0 [56.9, 63.0] & \relax 100.0 [99.6, 100.0] & \relax 100.0 [99.6, 100.0] \\
\relax Chi-square across grades, read once at year 10 & \relax 4.4 [3.3, 5.9] & \relax 66.0 [63.0, 68.9] & \relax 99.8 [99.3, 99.9] & \relax 100.0 [99.6, 100.0] & \relax 100.0 [99.6, 100.0] \\
\relax Chi-square across grades, read every year & \relax 20.8 [18.4, 23.4] & \relax 78.7 [76.1, 81.1] & \relax 99.8 [99.3, 99.9] & \relax 100.0 [99.6, 100.0] & \relax 100.0 [99.6, 100.0] \\
\bottomrule
\end{tabular}}
\end{table*}

